\section{Algoritmo de Thomas}

\begin{frame}
	\frametitle{\insertsection}

	Considere un sistema lineal tridiagonal expresado en la forma $Ax=d$
	\begin{equation}\label{eq:tridiagonal}
		\begin{bmatrix}
			b_{1}  & c_{1}  & 0      & \cdots  & 0       & 0       \\
			a_{2}  & b_{2}  & c_{2}  & 0       & \cdots  & \vdots  \\
			0      & a_{3}  & b_{3}  & c_{3}   & \ddots  & 0       \\
			0      & 0      & \ddots & \ddots  & \ddots  & 0       \\
			\vdots & \vdots & 0      & a_{n-1} & b_{n-1} & c_{n-1} \\
			0      & \cdots & 0      & 0       & a_{n}   & b_{n}
		\end{bmatrix}
		\begin{bmatrix}
			x_{1}   \\
			x_{2}   \\
			x_{3}   \\
			\vdots  \\
			x_{n-1} \\
			x_{n}
		\end{bmatrix}=
		\begin{bmatrix}
			d_{1}   \\
			d_{2}   \\
			d_{3}   \\
			\vdots  \\
			d_{n-1} \\
			d_{n}
		\end{bmatrix}.
	\end{equation}
	Este sistema lineal se convierte en forma triangular superior.
	$Ux=d^{\prime}$ donde la diagonal principal solo contiene unos (1's).
	Considere la fila 1 de la ecuación~\eqref{eq:tridiagonal}
	\begin{equation*}
		b_{1}x_{1}+
		c_{1}x_{2}=
		d_{1}.
	\end{equation*}
	Divida por $b_{1}$ para que la diagonal principal sea igual a $1$:
	\begin{equation*}
		x_{1}+
		\frac{c_{1}}{b_{1}}x_{2}=
		\frac{d_{1}}{b_{1}}.
	\end{equation*}
	Sea $c^{\prime}_{1}=\frac{c_{1}}{b_{1}}$ y
	$d^{\prime}_{1}=\frac{d_{1}}{b_{1}}$, entonces
	\begin{equation*}
		x_{1}+
		c^{\prime}_{1}x_{2}=
		d^{\prime}_{1}.
	\end{equation*}
	Considere la fila 2 de la ecuación~\eqref{eq:tridiagonal}
	\begin{equation*}
		a_{2}x_{1}+
		b_{2}x_{2}+
		c_{2}x_{3}=
		d_{2}.
	\end{equation*}
	Elimina $x_{1}$ usando $a_{2}$ multiplicado por la ecuación~\eqref{eq:tridiagonal}
	\begin{align*}
		a_{2}x_{1}+
		b_{2}x_{2}+
		c_{2}x_{3}-
		a_{2}\left(x_{1}+c^{\prime}_{1}x_{2}\right)            & =
		d_{2}-a_{2}d^{\prime}_{1}.                                 \\
		\left(b_{2}-a_{2}c^{\prime}_{1}\right)x_{2}+c_{2}x_{3} & =
		d_{2}-a_{2}d^{\prime}_{1}.
	\end{align*}
\end{frame}

\begin{frame}
	\frametitle{\insertsection}

	Divida por $b_{2}-a_{2}c^{\prime}_{1}$ para que la diagonal
	principal elemento es igual a $1$
	\begin{equation*}
		x_{2}+
		\frac{c_{2}}{b_{2}-a_{2}c^{\prime}_{1}}x_{3}=
		\frac{d_{2}-a_{2}d^{\prime}_{1}}{b_{2}-a_{2}c^{\prime}_{1}}.
	\end{equation*}
	Sea $c^{\prime}_{2}=\dfrac{c_{2}}{b_{2}-a_{2}c^{\prime}_{1}}$ y
	\begin{math}
		d^{\prime}_{2}=
		\dfrac{d_{2}-a_{2}d^{\prime}_{1}}{b_{2}-a_{2}c^{\prime}_{1}}
	\end{math},
	entonces
	\begin{equation*}
		x_{2}+
		c^{\prime}_{2}x_{3}=
		d^{\prime}_{2}.
	\end{equation*}
	Considere la fila 3 de la ecuación~\eqref{eq:tridiagonal}
	\begin{equation*}
		a_{3}x_{2}+
		b_{3}x_{3}+
		c_{3}x_{4}=
		d_{3}.
	\end{equation*}
	Elimina $x_{2}$ usando $a_{3}$ multiplicado por la
	ecuación~\eqref{eq:tridiagonal}
	\begin{align*}
		a_{3}x_{2}+
		b_{3}x_{3}+
		c_{3}x_{4}-
		a_{3}\left(x_{2}+c^{\prime}_{2}x_{3}\right) & =
		d_{3}-a_{3}d^{\prime}_{2}.                      \\
		\left(b_{3}-a_{3}c^{\prime}_{2}\right)x_{3}+
		c_{3}x_{4}                                  & =
		d_{3}-a_{3}d^{\prime}_{2}.
	\end{align*}
	Divida por $b_{3}-a_{3}c^{\prime}_{2}$ para que la diagonal principal
	elemento es igual a $1$
	\begin{equation*}
		x_{3}+
		\frac{c_{3}}{b_{3}-a_{3}c^{\prime}_{2}}x_{4}=
		\frac{d_{3}-a_{3}d^{\prime}_{2}}{b_{3}-a_{3}c^{\prime}_{2}}.
	\end{equation*}
	Sea $c^{\prime}_{3}=\dfrac{c_{3}}{b_{3}-a_{3}c^{\prime}_{2}}$ y
	\begin{math}
		d^{\prime}_{3}=
		\dfrac{d_{3}-a_{3}d^{\prime}_{2}}{b_{3}-a_{3}c^{\prime}_{2}}
	\end{math},
	entonces
	\begin{equation*}
		x_{3}+c^{\prime}_{3}x_{4}=
		d^{\prime}_{3}.
	\end{equation*}
\end{frame}

\begin{frame}
	\frametitle{\insertsection}

	Considere la $i$-ésima fila de la ecuación~\eqref{eq:tridiagonal}.
	Usando un método similar al que se muestra para las filas $2$ y $3$
	arriba, la $i$-ésima fila se puede escribir como
	\begin{equation}\label{eq:ith}
		x_{i}+
		c^{\prime}_{i}x_{i+1}=
		d^{\prime}_{i}.
	\end{equation}
	Donde
	\begin{align*}
		c^{\prime}_{i} & =
		\frac{c_{i}}{b_{i}-a_{i}c^{\prime}_{i-1}}. \\
		d^{\prime}_{i} & =
		\frac{d_{i}-a_{i}d^{\prime}_{i-1}}{b_{i}-a_{i}c^{\prime}_{i-1}}.
	\end{align*}
	Finalmente, considere la enésima fila de la ecuación~\eqref{eq:tridiagonal}
	\begin{equation*}
		a_{n}x_{n-1}+
		b_{n}x_{n}=
		d_{n}.
	\end{equation*}
	Elimina $x_{n-1}$ usando $a_{n}$ multiplicado por la ecuación~\eqref{eq:ith}.
	\begin{align*}
		a_{n}x_{n-1}+b_{n}x_{n}-a_{n}\left(x_{n-1}+c^{\prime}_{n-1}x_{n}\right) & =
		d_{n}-a_{n}d^{\prime}_{n-1}.                                                \\
		\left(b_{n}-a_{n}c^{\prime}_{n-1}\right)x_{n}                           & =
		d_{n}-a_{n}d^{\prime}_{n-1}.
	\end{align*}
	Divida por $b_{n}-a_{n}c^{\prime}_{n-1}$ para que el elemento de la
	diagonal principal sea igual a $1$.
	\begin{equation*}
		x_{n}=
		\frac{d_{n}-a_{n}d^{\prime}_{n-1}}{b_{n}-a_{n}c^{\prime}_{n-1}}.
	\end{equation*}
	En resumen, el sistema lineal tridiagonal~\eqref{eq:tridiagonal} se
	convierte en forma triangular superior $Ux=d^{\prime}$, donde $U$
	es una matriz tridiagonal con $1$'s en la diagonal principal,
	$c^{\prime}_{i}$'s en la diagonal superior y $0$'s en la diagonal
	inferior.
\end{frame}

\begin{frame}
	\frametitle{\insertsection}
	Los elementos de $U$ y $d^{\prime}$ se calculan usando las
	siguientes fórmulas de recurrencia:
	\begin{align*}
		c^{\prime}_{1} & =
		\frac{c_{1}}{b_{1}},\quad
		d^{\prime}_{1}=
		\frac{d_{1}}{b_{1}}, \\
		c^{\prime}_{i} & =
		\frac{c_{i}}{b_{i}-a_{i}c^{\prime}_{i-1}},\quad
		d^{\prime}_{i}=
		\frac{d_{i}-a_{i}d^{\prime}_{i-1}}{b_{i}-a_{i}c^{\prime}_{i-1}},\quad
		i=2,3,\dotsc,n.
	\end{align*}
	Una vez que el sistema está en forma triangular superior, podemos
	resolver $x$ usando sustitución hacia atrás:
	\begin{align*}
		x_{n} & =
		d^{\prime}_{n}, \\
		x_{i} & =
		d^{\prime}_{i}-
		c^{\prime}_{i}x_{i+1},\quad i=n-1,n-2,\dotsc,1.
	\end{align*}
	La ecuación matricial ahora se puede expresar en la forma
	$Ux=d^{\prime}$.
	\begin{equation}\label{eq:upperdiagonal}
		\begin{bmatrix}
			1      & c^{\prime}_{1} & 0              & \cdots         & 0      & 0                \\
			0      & 1              & c^{\prime}_{2} & 0              & \cdots & \vdots           \\
			0      & 0              & 1              & c^{\prime}_{3} & 0      & 0                \\
			0      & 0              & \ddots         & \ddots         & \ddots & 0                \\
			\vdots & \vdots         & 0              & 0              & 1      & c^{\prime}_{n-1} \\
			0      & \cdots         & 0              & 0              & 0      & 1
		\end{bmatrix}
		\begin{bmatrix}
			x_{1}   \\
			x_{2}   \\
			x_{3}   \\
			\vdots  \\
			x_{n-1} \\
			x_{n}
		\end{bmatrix}=
		\begin{bmatrix}
			d^{\prime}_{1}   \\
			d^{\prime}_{2}   \\
			d^{\prime}_{3}   \\
			\vdots           \\
			d^{\prime}_{n-1} \\
			d^{\prime}_{n}
		\end{bmatrix}.
	\end{equation}
	Una vez que el sistema está escrito en forma triangular superior,
	la solución de $x$ se puede encontrar mediante sustitución hacia
	atrás.
	Comenzando con la $n$-ésima fila de la ecuación~\eqref{eq:upperdiagonal}
	la solución para $x_{n}$ es $d^{\prime}_{n}$.
	\begin{equation*}
		x_{n}=d^{\prime}_{n}.
	\end{equation*}
	Considere la $i$-ésima fila de la ecuación~\eqref{eq:upperdiagonal}
	cuando $i=n-1,n-2,\dotsc,1$:
	\begin{equation*}
		x_{i}=
		d^{\prime}_{i}-
		c^{\prime}_{i}x_{i+1}.
	\end{equation*}
\end{frame}

\begin{frame}
	\frametitle{\insertsection}

	El algoritmo de Thomas se escribe de la siguiente manera.

	\begin{itemize}
		\item

		      Calcule los coeficientes modificados $c^{\prime}$ y
		      $d^{\prime}$ usando las fórmulas de recurrencia.
		      \begin{equation*}
			      \left\{
			      \begin{aligned}
				      c^{\prime}_{1} & =
				      \frac{c_{1}}{b_{1}},\quad
				      d^{\prime}_{1}=
				      \frac{d_{1}}{b_{1}}, \\
				      c^{\prime}_{i} & =
				      \frac{c_{i}}{b_{i}-a_{i}c^{\prime}_{i-1}},\quad
				      d^{\prime}_{i}=
				      \frac{d_{i}-a_{i}d^{\prime}_{i-1}}{b_{i}-a_{i}c^{\prime}_{i-1}},\quad
				      i=2,3,\dotsc,n.
			      \end{aligned}
			      \right.
		      \end{equation*}

		\item

		      Calcule la solución $x$ usando sustitución hacia atrás.
		      \begin{equation*}
			      \left\{
			      \begin{aligned}
				      x_{n} & =
				      d^{\prime}_{n}, \\
				      x_{i} & =
				      d^{\prime}_{i}-c^{\prime}_{i}x_{i+1},\quad i=n-1,n-2,\dotsc,1.
			      \end{aligned}
			      \right.
		      \end{equation*}
	\end{itemize}
\end{frame}
